\section{Sources, categories, and conventions}

Our primary manuscript source is \citet{AB_Grothendieck1981},
in the supplied scans and checked transcriptions. Archive scan
numbers and manuscript folios are different identifiers.
For example C3-399 is §48.X, folio500 verso; C3-489 is §48.XII,
folio547; C3-546 is §49.IV, folio576; C3-559--562 correspond to
folios583--584; C3-601--611 to folios604--609; C3-651--653 to
folios628--629; and C4-33--35 to §51.1, folios664--665.
An archive footer is not automatically a manuscript folio.

Unless a base is explicitly changed, schemes are over $\Z$.
An exact sequence of group schemes means an exact sequence of
fppf group sheaves; surjectivity on $R$-points is asserted
separately. Infinite inverse-limit projections are assigned their
actual topology. A constant group scheme $\underline A$ and a
roots-of-unity group scheme $\mu_n$ are not interchanged at
primes dividing $n$.

Fix

$$
r=\begin{pmatrix}0&1\\-1&1\end{pmatrix},\qquad
 s=\begin{pmatrix}0&-1\\1&0\end{pmatrix}.
$$

Thus $r^3=s^2=-I$. Put

$$
u_1(z)=\begin{pmatrix}1&0\\z&1\end{pmatrix},
 \qquad
 u_0(z)=\begin{pmatrix}1&-z\\0&1\end{pmatrix}.
$$

Then $rs=u_1(1)$ and $sr=-s^{-1}r=u_0(1)$.
C3-479 uses the opposite sign for the sigma matrix.
We keep the displayed convention and translate signs rather
than silently combining the two conventions.

For any commutative ring $R$ define

$$
\begin{split}
 T_r(R)&=\{aI+br:a^2+ab+b^2\in R^\times\},\\
 T_s(R)&=\{cI-ds:c^2+d^2\in R^\times\},\\
 N'_s(R)&=\left\{
 \begin{pmatrix}c&d\\-ed&ec\end{pmatrix}:
 e^2=1,\ c^2+d^2\in R^\times\right\}.
 \end{split}
$$

Write $n=a^2+ab+b^2$, $q=c^2+d^2$.
The prime in $N'_s$ specifies this signed normalizer, not an
unspecified full torus normalizer.

The principal matrix group in this paper is exactly

$$
M(R)=\{(x,y,f):x\in T_r(R),\ y\in N'_s(R),\
 f\in\GL_2(R),\ \det f=\det y=t\det x,\ t^2=1\}.
$$

The sign $e$ is determined by $ys=e\,sy$, and
$t=\det f/\det x$ is determined by the tuple. Equivalently,

$$
n=etq,\qquad \det f=eq,\qquad e^2=t^2=1.
$$

All projections and determinant functions used below are
morphisms of group schemes.

Put $M^{00}=\ker(e,t)$ and $M^+=\ker(\det f)$.
The seven-dimensional Borel fibre product called $SM$ below
is a different construction. The three-dimensional scheme
on C4-41 is also different from $M$.
For the lattice discussion, $E_L(\Gamma)$ is a finite-type
model, whereas

$$
\Env(\Gamma)=\varprojlim_L E_L(\Gamma)
$$

is the full affine inverse limit.

The manuscript's geometric unipotence condition means that,
after extension to an algebraic closure of the residue field,
the element is a value of a homomorphism from $\Ga$.
In characteristic $p$ this is stronger than merely being a
unipotent matrix. The word ``metabelian'' in §50 is read using
its local manuscript definition, finitely generated virtually
abelian. It is not replaced by the modern derived-length-two
condition.

The symbols $G_\Q$, $\widehat{\GT}$, and
$\widehat{\GT}_0$ always denote distinct objects. An assertion
about finite shadows of the first two GT relations is not an
assertion about the full GT group, and an assertion about
$G_\Q$ is not automatically an assertion about an unlabeled
``group of loop solutions.''

The supplied earlier AI returns are leads for investigation.
They are not cited as mathematical authorities. Literature
results used as theorems are identified by exact statements
and reductions; original calculations and counterexamples
are proved here. The readings register distinguishes a
visually inspected reading from an algebraically forced
correction and from an unresolved pen stroke.
